\providecommand{\weakconv}{\rightharpoonup}
\providecommand{\K}{\mathbb{K}}

\subsection{Замкнутый шар в гильбертовом пространстве слабо секвенциально компактен (теорема Банаха)}

\begin{definition}[Слабая секвенциальная компактность]
	Пусть \(E\) --- нормированное пространство. Множество
	\[
	M\subset E
	\]
	называется \textbf{слабо секвенциально компактным}, если из любой последовательности
	\[
	\{x_n\}\subset M
	\]
	можно выделить подпоследовательность \(\{x_{n_k}\}\), слабо сходящуюся к некоторому
	элементу \(x\in M\):
	\[
	x_{n_k}\weakconv x.
	\]
\end{definition}

\begin{remark}[Слабая сходимость в гильбертовом пространстве]
	Пусть \(H\) --- гильбертово пространство. По теореме Рисса--Фреше каждый функционал
	\(f\in H^*\) имеет вид
	\[
	f(z)=(z,y)
	\]
	для некоторого \(y\in H\). Поэтому
	\[
	x_n\weakconv x
	\quad\Longleftrightarrow\quad
	\forall y\in H
	\qquad
	(x_n,y)\to(x,y).
	\]
\end{remark}

\begin{theorem}[Банаха]
	Пусть \(H\) --- гильбертово пространство и
	\[
	\overline{B}_R(0)=\{x\in H:\|x\|\leqslant R\}.
	\]
	Тогда замкнутый шар \(\overline{B}_R(0)\) слабо секвенциально компактен.
	
	Иными словами, из любой последовательности
	\[
	x_n\in\overline{B}_R(0)
	\]
	можно выделить подпоследовательность \(\{x_{n_k}\}\), такую что
	\[
	x_{n_k}\weakconv x
	\]
	для некоторого
	\[
	x\in\overline{B}_R(0).
	\]
\end{theorem}

\begin{idea}
	Берём последовательность в шаре. Она может жить в несепарабельном пространстве,
	но сама эта последовательность порождает сепарабельное подпространство
	\[
	L=\overline{\operatorname{span}}\{x_1,x_2,\ldots\}.
	\]
	Внутри \(L\) диагональным методом добиваемся сходимости всех чисел
	\[
	(x_m,x_{n_k})
	\]
	при фиксированном \(m\). Это даёт сходимость скалярных произведений на плотном
	подпространстве.
	
	После этого строим функционал
	\[
	F(z)=\lim_{k\to\infty}(z,x_{n_k}),
	\]
	по теореме Рисса--Фреше предъявляем элемент \(x\in L\), и получаем слабую сходимость.
\end{idea}

\begin{proof}
	Пусть
	\[
	x_n\in\overline{B}_R(0),
	\qquad n\in\mathbb{N}.
	\]
	То есть
	\[
	\|x_n\|\leqslant R
	\qquad
	\forall n.
	\]
	Нужно выделить подпоследовательность, слабо сходящуюся к элементу того же шара.
	
	\textbf{1. Переход к подпространству, порождённому последовательностью.}
	
	Рассмотрим
	\[
	L=\overline{\operatorname{span}}\{x_1,x_2,\ldots\}\subset H.
	\]
	Тогда
	\[
	x_n\in L
	\qquad
	\forall n.
	\]
	Подпространство \(L\) замкнуто в \(H\), поэтому \(L\) само является гильбертовым пространством.
	Кроме того, множество
	\[
	L_0=\operatorname{span}\{x_1,x_2,\ldots\}
	\]
	плотно в \(L\).
	
	Достаточно сначала построить слабую сходимость в \(L\), а затем перейти ко всему \(H\)
	через ортогональное разложение
	\[
	H=L\oplus L^\perp.
	\]
	
	\textbf{2. Диагональный выбор подпоследовательности.}
	
	Для фиксированного \(m\in\mathbb{N}\) числовая последовательность
	\[
	(x_m,x_n),
	\qquad n=1,2,\ldots,
	\]
	ограничена, потому что по неравенству Коши--Буняковского
	\[
	|(x_m,x_n)|
	\leqslant
	\|x_m\|\cdot\|x_n\|
	\leqslant
	R^2.
	\]
	Значит, из неё можно выделить сходящуюся подпоследовательность.
	
	Теперь применим диагональный метод Кантора.
	
	\begin{itemize}
		\item Из \(\{x_n\}\) выделим подпоследовательность, на которой сходится
		\[
		(x_1,x_n).
		\]
		
		\item Из неё выделим подпоследовательность, на которой сходится
		\[
		(x_2,x_n).
		\]
		
		\item Продолжая процесс, получаем систему вложенных подпоследовательностей.
		
		\item Берём диагональную подпоследовательность и обозначаем её
		\[
		y_k=x_{n_k}.
		\]
	\end{itemize}
	
	Тогда для каждого фиксированного \(m\in\mathbb{N}\) существует предел
	\[
	\lim_{k\to\infty}(x_m,y_k).
	\]
	
	\textbf{3. Построение функционала на плотном подпространстве.}
	
	Покажем, что для любого
	\[
	z\in L_0=\operatorname{span}\{x_1,x_2,\ldots\}
	\]
	существует предел
	\[
	\lim_{k\to\infty}(z,y_k).
	\]
	Действительно, если
	\[
	z=\sum_{m=1}^{N}\alpha_m x_m,
	\]
	то
	\[
	(z,y_k)
	=
	\sum_{m=1}^{N}\alpha_m(x_m,y_k)
	\]
	с точностью до выбранного соглашения о линейности скалярного произведения. Каждое
	\((x_m,y_k)\) имеет предел при \(k\to\infty\), значит, предел имеет и конечная линейная
	комбинация.
	
	Зададим на \(L_0\) функционал
	\[
	F(z)=\lim_{k\to\infty}(z,y_k).
	\]
	Он линеен, поскольку предел сохраняет линейные операции.
	
	Кроме того, \(F\) ограничен. Для любого \(z\in L_0\)
	\[
	|F(z)|
	=
	\left|\lim_{k\to\infty}(z,y_k)\right|
	\leqslant
	\limsup_{k\to\infty}|(z,y_k)|.
	\]
	По неравенству Коши--Буняковского
	\[
	|(z,y_k)|
	\leqslant
	\|z\|\cdot\|y_k\|
	\leqslant
	R\|z\|.
	\]
	Следовательно,
	\[
	|F(z)|\leqslant R\|z\|.
	\]
	Значит, \(F\) ограничен на \(L_0\).
	
	Так как
	\[
	L=\overline{L_0},
	\]
	функционал \(F\) единственным образом продолжается до ограниченного линейного функционала
	на всём \(L\). Это продолжение снова обозначим через \(F\). При этом
	\[
	F\in L^*,
	\qquad
	\|F\|\leqslant R.
	\]
	
	\textbf{4. Предъявление слабого предела по теореме Рисса--Фреше.}
	
	Так как \(L\) --- гильбертово пространство, по теореме Рисса--Фреше существует элемент
	\(x\in L\), такой что
	\[
	F(z)=(z,x)
	\qquad
	\forall z\in L.
	\]
	Причём
	\[
	\|x\|=\|F\|\leqslant R.
	\]
	Следовательно,
	\[
	x\in\overline{B}_R(0).
	\]
	
	Теперь докажем, что
	\[
	y_k\weakconv x
	\qquad
	\text{в } L.
	\]
	По построению \(F\)
	\[
	F(z)=\lim_{k\to\infty}(z,y_k),
	\]
	а по выбору \(x\)
	\[
	F(z)=(z,x).
	\]
	Значит,
	\[
	(z,y_k)\to(z,x)
	\qquad
	\forall z\in L.
	\]
	По сопряжённой симметрии скалярного произведения это равносильно
	\[
	(y_k,z)\to(x,z)
	\qquad
	\forall z\in L.
	\]
	Следовательно,
	\[
	y_k\weakconv x
	\qquad
	\text{в } L.
	\]
	
	\textbf{5. Переход от \(L\) ко всему \(H\).}
	
	Так как \(L\) --- замкнутое подпространство гильбертова пространства, по теореме о проекции
	\[
	H=L\oplus L^\perp.
	\]
	Поэтому любой \(h\in H\) единственным образом представляется в виде
	\[
	h=l+l^\perp,
	\qquad
	l\in L,
	\qquad
	l^\perp\in L^\perp.
	\]
	Так как
	\[
	y_k\in L,
	\qquad
	x\in L,
	\]
	то
	\[
	(y_k,l^\perp)=0,
	\qquad
	(x,l^\perp)=0.
	\]
	Следовательно,
	\[
	(y_k,h)
	=
	(y_k,l+l^\perp)
	=
	(y_k,l)
	\to
	(x,l)
	=
	(x,l+l^\perp)
	=
	(x,h).
	\]
	То есть
	\[
	(y_k,h)\to(x,h)
	\qquad
	\forall h\in H.
	\]
	По критерию слабой сходимости в гильбертовом пространстве
	\[
	y_k\weakconv x
	\qquad
	\text{в } H.
	\]
	
	При этом уже доказано, что
	\[
	x\in\overline{B}_R(0).
	\]
	Значит, из любой последовательности в шаре \(\overline{B}_R(0)\) можно выделить
	подпоследовательность, слабо сходящуюся к элементу того же шара. Поэтому замкнутый шар
	в гильбертовом пространстве слабо секвенциально компактен.
\end{proof}
